This work establishes the EnKF as a dynamic and mechanistically grounded framework for knowledge transfer in upstream bioprocess modelling, allowing a system-specific mechanistic model to be progressively adapted to new conditions without reparametrisation. A central contribution lies in the dynamic evolution of parameter ensembles, which provide a temporally resolved view of how the model adapts to the new process. Unlike traditional approaches that treat model parameters as static, the EnKF framework tracks how parameter estimates evolve in response to real-time measurements. This time-resolved sensitivity analysis reveals which parameters are the most influential under new conditions, and when their impact on model behaviour becomes significant. Importantly, this is achieved without introducing artificial variability through parameter covariance inflation, thereby ensuring that changes in ensemble spread arise from real process data only.

Preserving mechanistic interpretability is particularly valuable for knowledge transfer, as it allows models to explain system behaviour beyond predictive performance. While hybrid and black-box approaches often prioritise generalisability, recent studies demonstrate that full mechanistic transparency can still be retained across diverse transfer contexts. Richelle \textit{et al.} demonstrated transfer from fed-batch to perfusion by simplifying model structure and introducing a ‘catch-all’ by-product variable to accommodate changes in operating mode \cite{Richelle2022Model-basedPerfusion}. Arndt \textit{et al.}  developed a scale-up workflow that quantified and compared model parameter distributions across bioreactor volumes, revealing when transfer is valid and when scale-dependent reparametrisation is necessary \cite{Arndt2021Model-basedSystems}. More recently, Lu \textit{et al.} incorporated time-series transcriptomics into a multi-cell-line learning (MCL) framework providing enzyme-level metabolic insight across CHO clones, though requiring high-performance computing and rich omics datasets \cite{Lu2024Multi-cell-lineModels}. This work complements these efforts by introducing time-resolved parameter trajectories, with the EnKF treating biologically meaningful kinetic parameters as dynamic quantities that evolve in response to process measurements. Compared with MCL, the presented method offers a computationally efficient alternative using only routine extracellular measurements, particularly suited for early process development where omics data may not be available.

While this method offers notable advantages in flexibility and interpretability, it relies on two key assumptions. First, the method  is designed to capture process-specific differences through parameter adaptation, without changing the underlying model structure. Second, the evolution of parameters is inherently conditioned on the initial ensemble drawn from the base system. In highly nonlinear systems, multiple parameter combinations may yield similar outputs, and the estimated trajectory represents one of the plausible solutions. These assumptions do not undermine the utility of the method but are essential to consider when interpreting biological meaning. Provided the model structure remains appropriate, the dynamic adaptation of parameters still offers valuable insight into system behaviour and sensitivity, especially when full reparametrisation is not feasible.

As the first demonstration EnKF dual state and parameter estimation in bioprocessing literatures, this work shows how dynamic model adaptation based on routine extracellular measurements can reveal system-specific behaviour without requiring full reparametrisation. The resulting dynamic parameter trajectories not only increase the flexibility of transferring mechanistic models across systems, but also provide biologically meaningful insight into process adaptation that extends beyond what is observable from concentration profiles alone. Beyond its methodological contributions, this framework offers a practical and interpretable route for reusing mechanistic models in new bioprocess settings where data are limited and timelines are constrained. From a regulatory perspective, the preservation of mechanistic structure and the transparent evolution of model parameters can support scientific justification for process comparability and risk assessment during tech transfer. In this context, the EnKF establishes itself as a robust and accessible tool for supporting knowledge-driven decision-making across development stages in upstream bioprocessing.